\documentclass[11pt]{amsart}
\usepackage{amsmath,amssymb,amsthm,booktabs}
\usepackage[margin=1.15in]{geometry}
\usepackage[hidelinks]{hyperref}
\newtheorem{theorem}{Theorem}[section]
\newtheorem{lemma}[theorem]{Lemma}
\newtheorem{proposition}[theorem]{Proposition}
\newtheorem{corollary}[theorem]{Corollary}
\newtheorem{conjecture}{Conjecture}
\theoremstyle{definition}
\newtheorem{definition}[theorem]{Definition}
\theoremstyle{remark}
\newtheorem{remark}[theorem]{Remark}
\newcommand{\PP}{\mathbb{P}}
\newcommand{\E}{\mathbb{E}}
\newcommand{\Z}{\mathbb{Z}}
\newcommand{\cL}{\mathcal{L}}
\newcommand{\doi}[1]{\href{https://doi.org/#1}{doi:\nolinkurl{#1}}}
\newcommand{\arxiv}[1]{\href{https://arxiv.org/abs/#1}{arXiv:\nolinkurl{#1}}}
\setlength{\emergencystretch}{2.5em}

\title[FA-1f above density one third]{Local relaxation of the one-dimensional Fredrickson--Andersen model for
infection density above one third}
\author{Yitzchak Shmalo}
\address{Einstein Institute of Mathematics, The Hebrew University of Jerusalem, Jerusalem, Israel}
\email{yspennstate@gmail.com}
\date{November 2026}
\subjclass[2020]{60K35, 82C22}
\keywords{kinetically constrained models, Fredrickson--Andersen model, relaxation to equilibrium, Lyapunov function}

\raggedbottom
\begin{document}
\begin{abstract}
In the one-spin facilitated Fredrickson--Andersen model on $\Z$, let $q$ be the equilibrium density of infected sites.
We show that for $q>1/3$ the law of the process seen in a fixed finite window converges to equilibrium exponentially
fast, up to an error equal to twice the probability that no infection starts within distance $vt$ of the origin, where
$v=v(q)>0$. Hence both conjectures of Martinelli, Shapira and Toninelli hold for $q>1/3$: the process converges from every initial law that has an infection almost surely, and it converges
exponentially fast when the law of the distance between the origin and the closest infection has an exponentially
decaying tail. Previously this was known for $q>1/2$. The new
ingredient is a Lyapunov function for a marked infection that follows the dynamics, corrected by the state of the site
behind it; the rest follows a block argument of Blondel, Cancrini, Martinelli, Roberto and
Toninelli.
\end{abstract}
\maketitle

\section{Introduction}\label{sec:intro}

We write configurations as $\eta\in\{0,1\}^{\Z}$, with $\eta_x=0$ meaning that $x$ is infected and $\eta_x=1$ that it
is healthy. For $q\in(0,1)$ we write $p=1-q$, and $\pi$ denotes the Bernoulli product measure under which each site is
infected with probability $q$. The one-spin facilitated Fredrickson--Andersen model (FA-1f) \cite{FA84} has generator
\begin{equation}\label{eq:gen}
Lg(\eta)=\sum_{x\in\Z}\bigl(1-\eta_{x-1}\eta_{x+1}\bigr)\bigl(\pi_x(g)-g\bigr)(\eta),
\end{equation}
where $\pi_x(g)$ is the conditional mean of $g$ given all spins except the one at $x$. In words, each site is resampled
from $\pi$ at rate one, but only while one of its two neighbours is infected. The measure $\pi$ is reversible. For a
probability measure $\nu$ on configurations, the initial law, $\nu P_t g$ denotes the mean of $g$ at time $t$. A local
function depends on the spins of a finite set of sites, its support. For a configuration $\eta$ and a site $x$ let
\[
R_x(\eta)=\min\{|x-y|:\ \eta_y=0\},\qquad R(\eta)=R_0(\eta),
\]
with $R_x=\infty$ when $\eta$ has no infection. Martinelli, Shapira and Toninelli \cite[Section~1.1]{MST25} state the
following two conjectures.

\begin{conjecture}\label{conj:1}
Let $q\in(0,1)$. If $\nu(R<\infty)=1$, then $\nu P_t f\to\pi(f)$ as $t\to\infty$ for every local $f$.
\end{conjecture}

\begin{conjecture}\label{conj:2}
Let $q\in(0,1)$. If $\nu(R>\rho)=O(e^{-\kappa\rho})$ as $\rho\to\infty$ for some $\kappa>0$, then for every local
$f$ the convergence $\nu P_t f\to\pi(f)$ is exponentially fast.
\end{conjecture}

The configuration with no infection is frozen, which is why some infection must be present. The model is not
attractive, and the constraint makes the dynamics degenerate, so neither monotone couplings nor the positive spectral
gap of FA-1f in equilibrium and in finite boxes \cite{CMRT08,CMRT09} give convergence from initial laws that are singular with respect to $\pi$.

For $q>1/2$, Blondel, Cancrini, Martinelli, Roberto and Toninelli establish exponential convergence
\cite[Theorem~2.1]{BCMRT13}, under the assumption that $\sup_x\nu(\theta^{R_x})<\infty$ for some $\theta>1$. A single infection violates
this assumption. Mountford and Valle \cite{MV19} treated a single infection, and initial laws as in
Conjecture~\ref{conj:2}, for $q$ close to one, and Ertul \cite{Ert21} proved cutoff on a finite interval from every
initial configuration for $q$ above a threshold that is roughly $0.76$. Exponential convergence for every $q$ is also known when the local
densities of the initial law with respect to $\pi$ grow slowly enough; this is due to Cancrini, Martinelli, Schonmann
and Toninelli \cite{CMST10}. In \cite{Shm26} the uniform assumption of \cite{BCMRT13} was replaced by one centred at the
origin, which gives both conjectures for $q>1/2$, from every initial law. For the general case
Gluchowski and Menz \cite{GM26}, who prove ergodicity for noisy variants of FA-1f and East, record Conjecture~\ref{conj:1}
as open.

\begin{theorem}\label{thm:main}
Fix $q\in(1/3,1)$ and a finite set $S\subset\Z$. Then there exist $v,c,C>0$, which depend on $q$ and $S$ alone, with
the following property: for each initial law $\nu$, each $f$ supported in $S$ and each time $t\ge0$,
\begin{equation}\label{eq:main}
\bigl|\nu P_tf-\pi(f)\bigr|\le\|f\|_\infty\bigl(Ce^{-ct}+2\,\nu(R>vt)\bigr).
\end{equation}
\end{theorem}

\begin{corollary}\label{cor:conj}
Conjectures~\ref{conj:1} and~\ref{conj:2} hold for every $q\in(1/3,1)$. Under the assumption of Conjecture~\ref{conj:2}
the rate of convergence can be taken to be $\min\{c,\kappa v\}$. In particular, started from one infection at any site,
the law of the spins in any finite window converges to its equilibrium law exponentially fast.
\end{corollary}

The term $\nu(R>vt)$ cannot be removed: when the only infection is at distance $\rho$ from $S$, the spins in $S$ cannot
change until a sequence of clock rings has crossed the gap, and this takes a time proportional to $\rho$
(Lemma~\ref{lem:speed}).

\emph{The method.} The proof follows the scheme of \cite[Section~3]{BCMRT13}. Near the end of the time interval one
works in a box of length of order $t$, cut into blocks of length proportional to $t$. As long as every block contains an
infection, the process agrees with a reversible chain whose spectral gap and logarithmic Sobolev constant are controlled by
a result of Cancrini, Martinelli, Roberto and Toninelli \cite[Proposition~3.4]{BCMRT13}, and that chain mixes. What has to be
supplied is a bound on the probability that some block is nearly free of infections, and for this \cite{BCMRT13} uses an
exponential moment of the distance $R_x$, which has a negative drift only for $q>1/2$ \cite[Proposition~4.1]{BCMRT13}.
Indeed, the nearest infection to $x$ moves towards $x$ at rate $q$ and away from it at rate up to $p$.

We follow instead one infection, the \emph{marked infection}, which moves towards $x$ whenever its neighbour on the side of
$x$ is infected and which, when it heals, steps to an infected neighbour (Definition~\ref{def:mark}). Its distance $D$
from $x$ still decreases at rate $q$ and increases at rate up to $p$. The gain comes from the site behind the mark: while that site is healthy, the mark
cannot step back, and while it is infected, it can heal in the next step. With $b$ the indicator that the site behind the mark
is infected, the function $W=D+ub$ with $u=p/(1+q)$ has drift at most $-q(3q-1)/(1+q)$ for every configuration
(Lemma~\ref{lem:drift}). This is negative exactly when $q>1/3$, and Remark~\ref{rem:third} shows that no choice of $u$ does
better. The same holds in a box whose exterior is frozen infected, which is the setting in which the block argument is run.

Table~\ref{tab:map} shows where each step is carried out.

\begin{table}[ht]
\caption{Organisation of the paper: for each section, the step of the proof it contains and the statements proved there.}\label{tab:map}
\begin{center}
\begin{tabular}{lll}
\toprule
Section & Content & Main statements\\
\midrule
\ref{sec:prelim} & graphical construction, box process, finite speed & Lemma~\ref{lem:speed}\\
\ref{sec:mark} & the marked infection and its drift & Lemmas~\ref{lem:drift}, \ref{lem:moment}\\
\ref{sec:box} & block argument in a box with infected exterior & Propositions~\ref{prop:gap}, \ref{prop:box}\\
\ref{sec:proof} & proof of the main results & Theorem~\ref{thm:main}, Corollary~\ref{cor:conj}\\
\ref{sec:remarks} & why the method stops at $q=1/3$ & Remark~\ref{rem:third}\\
\bottomrule
\end{tabular}
\end{center}
\end{table}

\section{Construction and finite speed}\label{sec:prelim}

This section fixes the graphical construction used throughout, defines the process in a box whose exterior is frozen
infected, and proves the finite speed estimate that lets us compare processes agreeing on an interval. Every constant
($c_0$, $c_1$, $C_1$, $C'$ and the others) is positive, is fixed where it is introduced, and keeps that value; no
constant changes from line to line.

\subsection{Clocks and spins}
Every site $x\in\Z$ carries a Poisson clock of rate one, and with its $j$-th ring we associate a spin $\omega_{x,j}$, equal
to $0$ with probability $q$ and to $1$ with probability $p$; all clocks and spins are independent. When the clock of $x$
rings for the $j$-th time, the spin at $x$ becomes $\omega_{x,j}$ provided a neighbour of $x$ is infected just before the
ring; otherwise nothing changes. Because the constraint involves only the two nearest-neighbour sites, this rule determines the evolution on $\Z$
almost surely for every starting configuration \cite{CMRT08}, and the result is the process with generator
\eqref{eq:gen}. We write $\PP_\eta$, $\E_\eta$ for its law and expectation from $\eta$, and $P_tg(\eta)=\E_\eta g(\eta(t))$.

\subsection{Boxes with infected exterior}
For an interval $\Lambda=\{\alpha,\dots,\beta\}$ of $\Z$ we use the process on $\{0,1\}^\Lambda$ in which the two exterior
neighbours $\alpha-1$ and $\beta+1$ count as infected forever and are never updated, and in which each site of $\Lambda$ follows the
rule above with the clocks and spins of the previous paragraph. We call it the \emph{box process} on $\Lambda$. It is a
continuous-time Markov chain on a finite set and is reversible with respect to the product measure $\pi_\Lambda$, the
restriction of $\pi$ to $\Lambda$. For $\xi\in\{0,1\}^\Lambda$ and $x\in\Lambda$ $R^\Lambda_x(\xi)$ denotes how far $x$
lies from the nearest site that is either infected in $\xi$ or equal to $\alpha-1$ or $\beta+1$. If $\eta$ agrees with
$\xi$ on $\Lambda$, then $R^\Lambda_x(\xi)\le R_x(\eta)$.

\subsection{Finite speed}
The next lemma is used twice to compare processes that agree on an interval and differ outside it.

\begin{lemma}\label{lem:speed}
Consider an interval $I\subset\Z$ and a finite set $S\subset I$ at distance $d\ge1$ from $\Z\setminus I$. Suppose that
$\eta(\cdot)$ and $\zeta(\cdot)$ are built from the same clocks and spins on a time interval $[\tau_0,\tau_0+w]$, such that
at every site of $I$ both follow the update rule above, while outside $I$ they evolve in an arbitrary way adapted to the
clocks and spins. If $\eta(\tau_0)$ and $\zeta(\tau_0)$ agree on $I$, then the probability that they differ at some site
of $S$ at some time of $[\tau_0,\tau_0+w]$ is at most $|S|(2w)^d/d!$, and at most $|S|4^{-d}$ when $d\ge8ew$.
\end{lemma}

\begin{proof}
Say that a site $y$ is in conflict at time $\tau$ if the two spins at $y$ differ at $\tau$. A site of $I$ that is not in
conflict at time $\tau_0$ can enter conflict only at a ring of its own clock. Both processes then use the same spin, so
the constraints differed just before the ring, and one of the two neighbours was in conflict at that moment. Suppose some
$y_0\in S$ is in conflict at a time of $[\tau_0,\tau_0+w]$. Take the ring at which $y_0$ last entered conflict and a
neighbour $y_1$ in conflict just before it; if $y_1\in I$, repeat with $y_1$ and that earlier time. This gives a
nearest-neighbour path $y_0,y_1,\dots$ and ring times of the clocks of $y_0,y_1,\dots$ that decrease strictly inside
$[\tau_0,\tau_0+w]$, and the path leaves $I$ after at least $d$ steps. Hence for some nearest-neighbour path
$y_0,\dots,y_{d-1}$ with $y_0\in S$ there are rings of the clocks of $y_0,\dots,y_{d-1}$ at decreasing times in
$[\tau_0,\tau_0+w]$. For a fixed path the expected number of such $d$-tuples of rings is $w^d/d!$ (this holds also if the
path revisits sites), and there are at most $|S|2^{d}$ paths. The bound follows from Markov's inequality, and the second
bound from $d!\ge(d/e)^d$.
\end{proof}

\section{The marked infection}\label{sec:mark}

Throughout this section $\Lambda=\{\alpha,\dots,\beta\}$ is an interval, the box process on $\Lambda$ is written $\xi(\cdot)$, and
$x\in\Lambda$ is fixed. The two \emph{frozen sites} $\alpha-1$ and $\beta+1$ count as infected. A site of $\Lambda\cup\{\alpha-1,\beta+1\}$
is called infected if it is frozen or infected in $\xi$.

\begin{definition}\label{def:mark}
A \emph{mark} is a site $y\in\Lambda\cup\{\alpha-1,\beta+1\}$ that is infected. For $y\ne x$ let $y^+=y+\operatorname{sign}(x-y)$
be the neighbour of $y$ on the side of $x$ and $y^-=y-\operatorname{sign}(x-y)$ the other neighbour; note that $y^+\in\Lambda$.
At each ring of a clock the box process is updated first, and the mark then moves according to the following rules; it
moves at no other time.
\begin{enumerate}
\item If $y\ne x$ and the clock of $y^+$ rings with spin $0$, the mark moves to $y^+$. (The update at $y^+$ is
allowed, since $y$ is infected, so $y^+$ is infected afterwards; it may have been infected before.)
\item If $y\in\Lambda$, $y\ne x$, and $y$ heals at a ring of its clock, the mark moves to $y^+$ if $y^+$ is infected and
to $y^-$ otherwise.
\item If $y=x$ and $x$ heals, the mark moves to $x-1$ if that site is infected and to $x+1$ otherwise.
\end{enumerate}
\end{definition}

A site heals only when one of its neighbours is infected, so in rules (2) and (3) the new position is infected, and in rule
(1) it is infected by construction. Hence the mark is always infected, and the pair $(\xi(\tau),y(\tau))$ is a Markov chain
on a finite set. A frozen mark moves only by rule (1). Let $D=|y-x|$ and let $b=1$ if $y\in\Lambda$, $y\ne x$ and $y^-$ is
infected (frozen or not), and $b=0$ otherwise; thus $b=0$ when $y=x$ and when the mark is frozen. Set
\begin{equation}\label{eq:consts}
u=\frac{p}{1+q},\qquad W=D+ub,\qquad \gamma=\frac{q(3q-1)}{1+q}.
\end{equation}
Write $\cL$ for the generator of the pair chain.

\begin{lemma}\label{lem:drift}
Let $q\in(1/3,1)$. Then $\gamma>0$, and $\cL W\le-\gamma$ at every state with $D\ge2$. Moreover, at every state, $D$ and
$b$ change only at rings of the clocks of $y^-$, $y$ and $y^+$ (of $x$ alone when $y=x$), and every change of $W$ has
absolute value at most $2$.
\end{lemma}

\begin{proof}
The last two statements are read off from Definition~\ref{def:mark}: the mark moves only at rings of the clocks of $y$ and
$y^+$, the value of $b$ depends on the spin at $y^-$ and on the position of the mark, and a move changes $D$ by one and $b$
by at most one, while $u<1$. Let $D\ge2$, so that $y^+\ne x$ and the site $y^+$ is not frozen. We consider four cases and
list the transitions that change $W$, with their rates and the changes of $W$; rings not listed leave $D$ and $b$
unchanged. We use the identities
\[
-q+2qu=-\gamma,\qquad -q+p(1-u)=-\gamma,
\]
which follow from $1-2u=(3q-1)/(1+q)$ and $1-u=2q/(1+q)$.

\emph{Case 1: $y\in\Lambda$, $y^-\in\Lambda$ and $y^-$ is healthy ($b=0$).} A ring of the clock of $y^+$ with spin $0$
(rate $q$) moves the mark to $y^+$, whose rear neighbour is now $y$, which is infected: the change is $-1+u$. A ring of
the clock of $y^-$ with spin $0$ (rate $q$) infects $y^-$, which is allowed since $y$ is infected: the change is $+u$.
The mark heals only if $y^+$ is infected, since $y^-$ is healthy; it then moves to $y^+$ and leaves the healthy site $y$
behind it, so the change is $-1$. The drift is at most $q(-1+u)+qu=-\gamma$.

\emph{Case 2: $y\in\Lambda$, $y^-\in\Lambda$ and $y^-$ is infected ($b=1$).} A ring of the clock of $y^+$ with spin $0$
(rate $q$) moves the mark forward with $y$ infected behind it: the change is $-1$. A ring of the clock of $y^-$ with spin
$1$ (rate $p$) heals $y^-$, which is allowed since $y$ is infected: the change is $-u$. A ring of the clock of $y$ with
spin $1$ (rate $p$) heals the mark. If $y^+$ is infected the mark moves forward and the change is $-1-u$; otherwise it
moves to $y^-$, its new value of $b$ is at most one, and the change is at most $+1$. The drift is at most
$-q-pu+p=-\gamma$.

\emph{Case 3: $y\in\Lambda$ and $y^-$ is frozen ($b=1$).} The forward move is as in Case 2, with change $-1$ at rate $q$.
The site $y^-$ never heals. If the mark heals and $y^+$ is healthy, it moves to the frozen site $y^-$, where $b=0$, so the
change is $1-u$; if $y^+$ is infected the change is $-1-u$. The drift is at most $-q+p(1-u)=-\gamma$.

\emph{Case 4: the mark is frozen ($b=0$).} The only move is rule (1), at rate $q$, after which the mark is at $y^+\in
\Lambda$ with the frozen site behind it, so $b=1$ and the change is $-1+u$. The drift is $-q(1-u)=-2q^2/(1+q)$, which is
at most $-\gamma$ because $2q^2\ge3q^2-q$ for $q\le1$.
\end{proof}

\begin{lemma}\label{lem:moment}
Let $q\in(1/3,1)$ and fix $a\in(0,1]$ with $12e^2a\le\gamma$. Put $\lambda=a\gamma/2$ and $K=3e^4+e^2$. For every
initial state of the pair chain and every $\tau\ge0$,
\begin{equation}\label{eq:moment}
\E\,e^{aR^\Lambda_x(\xi(\tau))}\le\E\,e^{aW(\tau)}\le e^{-\lambda\tau}\,\E\,e^{aW(0)}+\frac K\lambda .
\end{equation}
\end{lemma}

\begin{proof}
The first inequality holds because the mark is infected, so $R^\Lambda_x\le D\le W$. Let $g=e^{aW}$. For $|z|\le2$ we have
$e^z-1\le z+\tfrac12z^2e^{|z|}$, so a change $\Delta$ of $W$ changes $g$ by at most $g\,(a\Delta+2e^2a^2)$, since $a\le1$.
By Lemma~\ref{lem:drift} the transitions that change $W$ have total rate at most three. At a state with $D\ge2$ this gives
\[
\cL g\le g\,\bigl(a\,\cL W+6e^2a^2\bigr)\le g\,\Bigl(-a\gamma+\frac{a\gamma}2\Bigr)=-\lambda g.
\]
At a state with $D\le1$ we have $W\le2$ and every change of $W$ is at most $2$, so $\cL g\le3e^{4a}\le3e^4$, while
$\lambda g\le e^2$ because $\lambda\le1/2$ and $g\le e^{2a}\le e^2$. Hence $\cL g\le-\lambda g+K$ at every state. The pair chain has finitely many
states, so $\frac{d}{d\tau}\E g(\tau)=\E\cL g(\tau)\le-\lambda\E g(\tau)+K$, and the second inequality follows.
\end{proof}

\section{Relaxation in a box}\label{sec:box}

This section is the block argument of \cite[Section~3]{BCMRT13}, written for an initial law given through the probability
of nearly empty blocks; we include the proof because the estimate on leaving the good set is used in a form slightly
different from \cite[Proposition~3.1]{BCMRT13}. Let $\Lambda$ be an interval partitioned into $k$ consecutive intervals
$\Lambda_1,\dots,\Lambda_k$, the blocks, each of length between $\ell$ and $M$. Let
\[
B=\{\xi:\ \text{each block contains an infection of }\xi\}
\]
and
\[
B_2=\{\xi:\ \text{each block contains at least two infections of }\xi\}.
\]

\begin{proposition}[{\cite[Proposition~3.4]{BCMRT13}, \cite[Theorem~6.4]{CMRT09}}]\label{prop:gap}
Let $q\in(0,1)$. There is $c_0=c_0(q)>0$ with the following property. Let $A$ be an interval, let $\Omega_A^*$ be the set
of configurations on $A$ with at least one infection, and let $\pi_A^*=\pi_A(\cdot\mid\Omega_A^*)$. The chain on
$\Omega_A^*$ in which each site of $A$ is resampled from $\pi$ at rate one while it has an infected neighbour in $A$ (the
exterior counting as healthy), has spectral gap at least $c_0$ and logarithmic Sobolev constant at most $|A|/c_0$ with
respect to $\pi_A^*$.
\end{proposition}

Note that $\Omega_A^*$ is closed for this chain: the last infection of $A$ has no infected neighbour and cannot heal.

\begin{proposition}\label{prop:box}
Let $q\in(0,1)$. There is $c_1=c_1(q)>0$ such that the following holds. Let $\Lambda$, the blocks $\Lambda_1,\dots,\Lambda_k$ and the sets $B$, $B_2$ be as at the start of this section, with
$kp^\ell\le1/2$, let $f$ have support in $\Lambda$, let $h>0$, and let $\xi(\cdot)$ be the box process on $\Lambda$ with
initial law $\mu$. Put $\varepsilon=\sup_{0\le w\le h}\PP_\mu(\xi(w)\notin B_2)$. Then
\[
\bigl|\E_\mu f(\xi(h))-\pi(f)\bigr|\le\|f\|_\infty\Bigl(2(1+kh)\varepsilon+4kp^\ell+
2\exp\Bigl\{-\frac{c_0h}2+|\Lambda|\log\frac1{\min(p,q)}\,e^{-c_1h/(2M)}\Bigr\}\Bigr).
\]
\end{proposition}

\begin{proof}
\emph{The chain on $B$.} Let $\hat\xi(\cdot)$ be the chain on $B$ that uses the clocks and spins of the box process and
performs each update of the box process except those that would leave $B$, which it ignores. The box process satisfies
detailed balance with respect to $\pi_\Lambda$, because the constraint at a site does not depend on the spin at that site.
Removing the transitions out of $B$ keeps detailed balance on $B$, so $\hat\xi$ is reversible with respect to
$\hat\pi=\pi_\Lambda(\cdot\mid B)$. This measure is the product over the blocks of $\pi^*_{\Lambda_i}$, in the notation of
Proposition~\ref{prop:gap}.

Consider on $B$ the product chain in which each block evolves independently as the chain of Proposition~\ref{prop:gap} on
$\Omega^*_{\Lambda_i}$. Every transition of the product chain is a transition of $\hat\xi$ with the same rate: if a site
has an infected neighbour within its block, it has one in $\Lambda\cup\{\alpha-1,\beta+1\}$, and the result stays in $B$. Both
chains are reversible with respect to $\hat\pi$, so the Dirichlet form of $\hat\xi$ dominates that of the product chain.
Spectral gaps and logarithmic Sobolev constants tensorize \cite{DSC96}, so the product chain, and hence $\hat\xi$, has
gap at least $c_0$ and logarithmic Sobolev constant at most $M/c_0$. By Gross's theorem \cite[Theorem~3.5]{DSC96} there is
$c_1=c_1(q)>0$ such that the semigroup $\hat P$ of $\hat\xi$ satisfies $\|\hat P_\tau G\|_{L^r(\hat\pi)}\le\|G\|_{L^2(\hat\pi)}$
for $r=1+e^{c_1\tau/M}$.

\emph{Mixing of $\hat\xi$.} Let $F=f-\hat\pi(f)$, so that $\|F\|_{L^2(\hat\pi)}\le2\|f\|_\infty$, and let $r=1+e^{c_1h/(2M)}$.
Every state of $B$ has $\hat\pi$-probability at least $\min(p,q)^{|\Lambda|}$, and for any function $G$ on $B$ and any
state $\sigma$, $|G(\sigma)|\le\hat\pi(\sigma)^{-1/r}\|G\|_{L^r(\hat\pi)}$. Using the gap on $[0,h/2]$ and hypercontractivity
on $[h/2,h]$,
\begin{align*}
\|\hat P_hF\|_\infty&\le\min(p,q)^{-|\Lambda|/r}\,\|\hat P_{h/2}(\hat P_{h/2}F)\|_{L^r(\hat\pi)}
\le\min(p,q)^{-|\Lambda|/r}\,\|\hat P_{h/2}F\|_{L^2(\hat\pi)}\\
&\le2\|f\|_\infty\exp\Bigl\{|\Lambda|\log\frac1{\min(p,q)}\,e^{-c_1h/(2M)}-\frac{c_0h}2\Bigr\}.
\end{align*}

\emph{Coupling.} Start $\hat\xi$ from $\xi(0)$ if $\xi(0)\in B$ and from the configuration with every site infected
otherwise, and run both chains with the same clocks and spins. They agree up to time $h$ unless $\xi(0)\notin B$, or some
ring at a time $w\le h$ would take $\xi(w-)=\hat\xi(w-)$ out of $B$. Such a ring heals a site that is the only infection
of its block, so $\xi(w-)\notin B_2$, and the ring belongs to the clock of one of at most $k$ sites. The rings of each clock
form a Poisson process of rate one and the event that its site is the only infection of its block just before time $w$ is
predictable, so the expected number of such rings in $[0,h]$ is at most $\int_0^hk\,\PP_\mu(\xi(w)\notin B_2)\,dw\le kh
\varepsilon$. Since $B_2\subset B$ and the supremum defining $\varepsilon$ includes $w=0$, also
$\PP_\mu(\xi(0)\notin B)\le\varepsilon$. Hence the two chains differ at time $h$ with probability at most $(1+kh)\varepsilon$, and
\[
\bigl|\E_\mu f(\xi(h))-\E\,\hat P_hf(\hat\xi(0))\bigr|\le2\|f\|_\infty(1+kh)\varepsilon .
\]

\emph{Conditioning.} A block of length at least $\ell$ is free of infections with $\pi$-probability at most $p^\ell$, so
$\pi_\Lambda(B^c)\le kp^\ell\le1/2$. Since the support of $f$ lies in $\Lambda$,
\[
\bigl|\hat\pi(f)-\pi(f)\bigr|=\frac{|\pi_\Lambda((f-\pi(f))\mathbf 1_{B^c})|}{\pi_\Lambda(B)}\le4\|f\|_\infty\,kp^\ell ,
\]
where $\mathbf 1_{B^c}$ is the indicator function of $B^c$.
Writing $\hat P_hf=\hat\pi(f)+\hat P_hF$ and collecting the three estimates proves the proposition.
\end{proof}

\section{Proofs of the main results}\label{sec:proof}

Throughout this section $q\in(1/3,1)$ is fixed, and $m\ge0$ is an integer such that $S\subset[-m,m]$. Take $u,\gamma$ from \eqref{eq:consts}, $a,\lambda,K$ from
Lemma~\ref{lem:moment}, and $c_0,c_1$ from Propositions~\ref{prop:gap} and~\ref{prop:box}. Set
\[
\delta=\min\Bigl\{\frac12,\frac\lambda{64ea}\Bigr\},\qquad v=\frac\lambda{8a}=\frac\gamma{16},
\]
so that $a(8e\delta+v)\le\lambda/4$ and $\lambda(1-\delta)\ge\lambda/2$. Choose $\varepsilon_0\in(0,1)$ such that
\begin{equation}\label{eq:eps}
50\log\frac1{\min(p,q)}\,e^{-c_1/(4\varepsilon_0)}\le\frac{c_0}4 .
\end{equation}
For $t>0$ let
\[
h=\delta t,\quad s=t-h,\quad n=m+\lceil8eh\rceil,\quad \Lambda=\{-n,\dots,n\},\quad N=n+\lceil8et\rceil,\quad
\Lambda'=\{-N,\dots,N\}.
\]
Divide $\Lambda$ into $k=\lceil|\Lambda|/(\varepsilon_0h)\rceil$ consecutive blocks of lengths $\lfloor|\Lambda|/k\rfloor$
and $\lceil|\Lambda|/k\rceil$, and let $\ell$ and $M$ be these two numbers. Since $16eh\le|\Lambda|\le2m+16eh+3$, there
is $h_0=h_0(q,m)$ such that for $h\ge h_0$ we have $|\Lambda|\le50h$, $k\le1+50/\varepsilon_0$,
$\ell\ge\varepsilon_0h/2$, $M\le2\varepsilon_0h$ and $kp^\ell\le1/2$. All later constants are functions of $q$ and $m$ only.

\begin{proof}[Proof of Theorem~\ref{thm:main}]
Take $t$ with $\delta t\ge h_0$ and a configuration $\eta$ such that $R(\eta)\le vt$, and pick an infected site $y_0$
of $\eta$ with $|y_0|\le vt$. Let $\eta(\cdot)$ be FA-1f on $\Z$ from $\eta$. Let $\zeta(\cdot)$ be the box process on
$\Lambda'$ started from the restriction of $\eta$, and let $\xi(\cdot)$, on the time interval $[s,t]$, be the box process
on $\Lambda$ started at time $s$ from the restriction of $\zeta(s)$; all three use the same clocks and spins.

\emph{Step 1: two comparisons.} The processes $\eta$ and $\zeta$ follow the update rule at the sites of $\Lambda'$ and
agree there at time $0$, and the distance from $S$ to $\Z\setminus\Lambda'$ is at least $N-m\ge8et$. The processes
$\zeta$ and $\xi$ follow the rule at the sites of $\Lambda$ on $[s,t]$ and agree there at time $s$, and the distance from
$S$ to $\Z\setminus\Lambda$ is at least $n-m\ge8eh$. By Lemma~\ref{lem:speed},
\begin{equation}\label{eq:step1}
\bigl|\E_\eta f(\eta(t))-\E f(\xi(t))\bigr|\le2\|f\|_\infty|S|\bigl(4^{-8et}+4^{-8eh}\bigr).
\end{equation}

\emph{Step 2: a moment bound at time $s$.} Let $x\in\Lambda$. Since $v<1\le8e$, the site $y_0$ lies in $\Lambda'$. Run the
pair chain of Section~\ref{sec:mark} in the box $\Lambda'$, for the site $x$, with the mark started at $y_0$. Then
$W(0)\le|x-y_0|+1\le m+2+(8e\delta+v)t$, and Lemma~\ref{lem:moment} gives
\[
\E\,e^{aR^{\Lambda'}_x(\zeta(s))}\le\exp\bigl\{a(m+2)+a(8e\delta+v)t-\lambda(1-\delta)t\bigr\}+\frac K\lambda
\le e^{a(m+2)}+\frac K\lambda=:K_2 .
\]

\emph{Step 3: a moment bound on $[s,t]$.} Given the history up to time $s$, run the pair chain in the box $\Lambda$, for
the site $x$, with the mark started at time $s$ at an infected site of $\Lambda\cup\{\pm(n+1)\}$ nearest to $x$. Since
$\xi(s)$ agrees with $\zeta(s)$ on $\Lambda$ and the sites $\pm(n+1)$ count as infected for $\xi$, its initial value of
$W$ is at most $R^\Lambda_x(\xi(s))+1\le R^{\Lambda'}_x(\zeta(s))+1$. By the Markov property at time $s$, Lemma~\ref{lem:moment}
and Step 2, for $0\le w\le h$,
\[
\E\,e^{aR^\Lambda_x(\xi(s+w))}\le e^aK_2+\frac K\lambda=:K_3,\qquad\text{so}\qquad
\PP\bigl(R^\Lambda_x(\xi(s+w))\ge j\bigr)\le K_3e^{-aj}\quad(j\ge0).
\]

\emph{Step 4: nearly empty blocks.} Cut each block into a left and a right half of lengths at least $\lfloor\ell/2\rfloor$,
and let $x$ run over the $2k$ sites, one in the middle of each half. A block with at most one infection has a half with no infection;
the middle site of that half is at distance at least $\ell/4-1$ from every infected site, frozen sites included, because all of
these lie outside the half. By Step 3 and $\ell\ge\varepsilon_0h/2$, for $0\le w\le h$,
\[
\PP\bigl(\xi(s+w)\notin B_2\bigr)\le2kK_3e^{a}e^{-a\varepsilon_0h/8}.
\]

\emph{Step 5: conclusion.} Apply Proposition~\ref{prop:box} to $\xi$ on $[s,t]$, with $\mu$ the law of $\xi(s)$. Since
$h/(2M)\ge1/(4\varepsilon_0)$ and $|\Lambda|\le50h$, condition \eqref{eq:eps} bounds the last term there by
$2\|f\|_\infty e^{-c_0h/4}$, and with Step 4,
\[
\bigl|\E f(\xi(t))-\pi(f)\bigr|\le\|f\|_\infty\Bigl(4(1+kh)kK_3e^ae^{-a\varepsilon_0h/8}+4kp^{\varepsilon_0h/2}
+2e^{-c_0h/4}\Bigr).
\]
Together with \eqref{eq:step1} and $h=\delta t$ this gives $c,C_1>0$ with $|P_tf(\eta)-\pi(f)|\le C_1\|f\|_\infty e^{-ct}$
for all $t\ge h_0/\delta$ and all $\eta$ with $R(\eta)\le vt$. Put $C=\max\{C_1,2e^{ch_0/\delta}\}$; since
$|P_tf-\pi(f)|\le2\|f\|_\infty$ always,
\begin{equation}\label{eq:pointwise}
\sup_{\eta:\ R(\eta)\le vt}\bigl|P_tf(\eta)-\pi(f)\bigr|\le C\|f\|_\infty e^{-ct}\qquad(t\ge0).
\end{equation}
Integrate this bound against $\nu$ over $\{R\le vt\}$ and use $|P_tf-\pi(f)|\le2\|f\|_\infty$ on $\{R>vt\}$; the
result is \eqref{eq:main}.
\end{proof}

\begin{proof}[Proof of Corollary~\ref{cor:conj}]
When $\nu(R<\infty)=1$, the term $\nu(R>vt)$ tends to zero, and Conjecture~\ref{conj:1} follows from \eqref{eq:main}. Under
the assumption of Conjecture~\ref{conj:2}, $\nu(R>vt)\le C'e^{-\kappa vt}$ for some $C'$ and all $t\ge0$, which gives the rate
$\min\{c,\kappa v\}$. For one infection at a site $x_0$, $R=|x_0|$ and the second term in \eqref{eq:main} is zero once
$t\ge|x_0|/v$.
\end{proof}

\section{The threshold one third}\label{sec:remarks}

\begin{remark}\label{rem:third}
Among functions of the form $D+ub$, the constant $u=p/(1+q)$ is the best one, and none has negative drift for $q\le1/3$.
Indeed, in Case~1 of the proof of Lemma~\ref{lem:drift} with $y^+$ healthy, the drift of $D+ub$ equals $-q+2qu$, which is
negative only if $u<1/2$. In Case~2 with $y^+$ healthy and $y^{--}=y^--\operatorname{sign}(x-y)$ infected, the mark
retreats onto a site whose rear neighbour is infected, and the drift equals $-q-pu+p$, which is negative only if
$u>(1-2q)/p$. Both conditions hold for some $u$ if and only if $(1-2q)/p<1/2$, that is $q>1/3$, and the value
$u=p/(1+q)$ makes the two drifts equal. This is a limitation of the test function, not a property of the process: for
$q\le1/3$ it is not known whether the two conjectures hold, even from a single infection. A function that also records the
sites two steps behind the mark, or a mark that is allowed to jump to a more distant infection, might lower the threshold.
\end{remark}

\section*{Acknowledgments and funding}

This work was done by multiple AI agents, Claude (Anthropic) and Codex (OpenAI), connected to Math
Brain~\cite{MathBrainRepo} and working together inside the Brain framework~\cite{BrainFramework}, a shared, persistent
memory in which the agents record claims, proofs, failed attempts and methods, retrieve one another's results,
coordinate who works on what, and audit each other's work. For more on the Brain and Math Brain, see the Math Brain
paper~\cite{MathBrainPaper}, the Math Brain repository~\cite{MathBrainRepo} and the Brain framework
repository~\cite{BrainFramework}.

The research presented in this paper was supported by the European Research Council (ERC) under the European
Union's Horizon Europe research and innovation programme (grant agreement No.\ 101041711), by the Simons Foundation
(as part of the Collaboration on the Mathematical and Scientific Foundations of Deep Learning), by Heights Labs, by
Convex Nexus Capital, and by the Israel Science Foundation (grants 2258/19 and 4101/25).

\section*{Code availability}

The Python script in the directory \texttt{checks} of the paper's repository computes, in exact rational arithmetic, the
drift of $W$ at every local configuration around the mark used in the proof of Lemma~\ref{lem:drift}, for $q=j/300$ with
$100<j<300$, and checks that its maximum over configurations with $D\ge2$ equals $-\gamma$. It also checks the two
drifts of Remark~\ref{rem:third}.

\section*{Declaration of competing interest}

The author holds equity in Heights Labs and Convex Nexus Capital and has a management role at Convex Nexus Capital.

\section*{Declaration of generative AI and AI-assisted technologies}

This paper is part of a broader research project on proving open problems with AI systems. The argument was found by an
OpenAI Codex agent, checked by Anthropic Claude agents, which also wrote the computer check and drafted the text. The
conjectures are due to Martinelli, Shapira and Toninelli \cite{MST25}, the block argument of Section~\ref{sec:box} is that
of Blondel, Cancrini, Martinelli, Roberto and Toninelli \cite{BCMRT13}, and the spectral gap bound used there is due to
Cancrini, Martinelli, Roberto and Toninelli \cite{CMRT09}. The author directed the work and is responsible for the content
of the paper and for decisions concerning its dissemination.

\appendix
\section{Notation}\label{app:notation}

\begin{center}
\small
\begin{tabular}{p{3.4cm}p{7.4cm}p{3.2cm}}
\toprule
Symbol & Meaning & First appears\\
\midrule
$\eta$, $\eta_x$ & a configuration on $\Z$ and its spin at $x$ ($0$ infected) & Section~\ref{sec:intro}\\
$q$, $p$ & density of infected sites under $\pi$, and $1-q$ & Section~\ref{sec:intro}\\
$\pi$, $\pi_x$ & product measure, conditional mean at $x$ & \eqref{eq:gen}\\
$L$, $P_t$ & generator and semigroup of FA-1f & \eqref{eq:gen}\\
$g$ & a function on configurations & \eqref{eq:gen}\\
$\nu$ & initial law & Section~\ref{sec:intro}\\
$\rho$, $\kappa$, $\theta$ & parameters in the conjectures and in prior results & Section~\ref{sec:intro}\\
$\omega_{x,j}$ & spin carried by the $j$-th ring of the clock of $x$ & Section~\ref{sec:prelim}\\
$\PP_\eta$, $\E_\eta$ & law and expectation of FA-1f from $\eta$ & Section~\ref{sec:prelim}\\
$\alpha$, $\beta$, $\xi$ & endpoints of $\Lambda$, a configuration of the box process & Section~\ref{sec:prelim}\\
$\zeta$, $\tau$, $\tau_0$ & a second process, time, initial time & Lemma~\ref{lem:speed}\\
$R_x$, $R$ & distance from $x$, from $0$, to the nearest infection & Section~\ref{sec:intro}\\
$S$, $f$ & finite set, function with support in $S$ & Theorem~\ref{thm:main}\\
$v$, $c$, $C$ & constants of Theorem~\ref{thm:main} & Theorem~\ref{thm:main}\\
$\Lambda$, $\pi_\Lambda$, $R^\Lambda_x$ & interval, restricted product measure, distance in the box & Section~\ref{sec:prelim}\\
$I$, $d$, $w$ & interval, distance from $S$ to its complement, time span & Lemma~\ref{lem:speed}\\
$y$, $y^\pm$, $D$, $b$ & mark, its neighbours, its distance to $x$, rear indicator & Definition~\ref{def:mark}\\
$u$, $W$, $\gamma$ & weight, Lyapunov function, drift bound & \eqref{eq:consts}\\
$\cL$ & generator of the pair chain & Section~\ref{sec:mark}\\
$a$, $\lambda$, $K$ & constants of the moment bound & Lemma~\ref{lem:moment}\\
$z$, $\Delta$ & real variable, a change of $W$ & proof of Lemma~\ref{lem:moment}\\
$A$, $\Omega^*_A$, $\pi^*_A$ & interval, configurations with an infection, conditioned measure & Proposition~\ref{prop:gap}\\
$F$, $G$, $r$, $\sigma$ & centred function, test function, exponent, state of $\hat\xi$ & proof of Proposition~\ref{prop:box}\\
$\Lambda_i$, $k$, $\ell$, $M$ & blocks, their number, minimal and maximal length & Section~\ref{sec:box}\\
$B$, $B_2$ & at least one, at least two infections per block & Section~\ref{sec:box}\\
$c_0$, $c_1$ & gap and hypercontractivity constants & Propositions~\ref{prop:gap}, \ref{prop:box}\\
$\hat\xi$, $\hat\pi$, $\hat P$ & chain on $B$, its reversible measure, its semigroup & Proposition~\ref{prop:box}\\
$\delta$, $\varepsilon_0$, $h$, $s$, $n$, $\Lambda'$ & parameters of the proof & Section~\ref{sec:proof}\\
$K_2$, $K_3$, $h_0$, $C_1$, $C'$ & constants of the proof & Section~\ref{sec:proof}\\
$m$, $N$, $\mu$ & half-width of a window containing $S$, half-width of $\Lambda'$, initial law in a box & Sections~\ref{sec:box}, \ref{sec:proof}\\
$x_0$ & site of a single initial infection & proof of Corollary~\ref{cor:conj}\\
\bottomrule
\end{tabular}
\end{center}

\begin{thebibliography}{99}
\bibitem{BCMRT13} O.~Blondel, N.~Cancrini, F.~Martinelli, C.~Roberto and C.~Toninelli, Fredrickson--Andersen one spin
facilitated model out of equilibrium, \emph{Markov Process.\ Related Fields} \textbf{19} (2013), no.~3, 383--406.
\arxiv{1205.4584}
\bibitem{CMRT08} N.~Cancrini, F.~Martinelli, C.~Roberto and C.~Toninelli, Kinetically constrained spin models,
\emph{Probab.\ Theory Related Fields} \textbf{140} (2008), no.~3--4, 459--504. \doi{10.1007/s00440-007-0072-3}
\bibitem{CMRT09} N.~Cancrini, F.~Martinelli, C.~Roberto and C.~Toninelli, Facilitated spin models: recent and new
results, in \emph{Methods of Contemporary Mathematical Statistical Physics}, Lecture Notes in Math.\ \textbf{1970},
Springer, Berlin, 2009, 307--340.
\bibitem{CMST10} N.~Cancrini, F.~Martinelli, R.~Schonmann and C.~Toninelli, Facilitated oriented spin models: some
non equilibrium results, \emph{J.\ Stat.\ Phys.} \textbf{138} (2010), no.~6, 1109--1123.
\doi{10.1007/s10955-010-9923-x}
\bibitem{DSC96} P.~Diaconis and L.~Saloff-Coste, Logarithmic Sobolev inequalities for finite Markov chains, \emph{Ann.\
Appl.\ Probab.} \textbf{6} (1996), no.~3, 695--750.
\bibitem{Ert21} A.~Ertul, Cutoff for the Fredrickson--Andersen one spin facilitated model, preprint, 2021.
\arxiv{2103.00019}
\bibitem{FA84} G.~H.~Fredrickson and H.~C.~Andersen, Kinetic Ising model of the glass transition, \emph{Phys.\ Rev.\
Lett.} \textbf{53} (1984), no.~13, 1244--1247. \doi{10.1103/PhysRevLett.53.1244}
\bibitem{GM26} M.~Gluchowski and G.~Menz, Ergodicity without monotonicity: patch-positive spin systems, preprint, 2026.
\arxiv{2611.17861}
\bibitem{MST25} F.~Martinelli, A.~Shapira and C.~Toninelli, Long time behaviour of one facilitated kinetically
constrained models: results and open problems, preprint, 2025. \arxiv{2510.20461}
\bibitem{MV19} T.~Mountford and G.~Valle, Exponential convergence for the Fredrickson--Andersen one-spin facilitated
model, \emph{J.\ Theoret.\ Probab.} \textbf{32} (2019), no.~1, 282--302. \doi{10.1007/s10959-018-0812-6}
\bibitem{Shm26} Y.~Shmalo, Relaxation of the one-dimensional Fredrickson--Andersen model from a single infection at high
density, preprint, 2026. \doi{10.5281/zenodo.24294310}
\bibitem{MathBrainRepo} Y.~Shmalo, \emph{Math Brain: public snapshot, code and tests}, 2026.
\url{https://github.com/yspennstate/math-brain-public/tree/5427734deee9be643d6dfbdd5e7d4baeeefe4a4a}
\bibitem{MathBrainPaper} Y.~Shmalo, \emph{Math Brain: shared, auditable memory for multi-agent mathematical research},
October 2026.
\url{https://github.com/yspennstate/math-brain-public/tree/5427734deee9be643d6dfbdd5e7d4baeeefe4a4a/papers/math-brain-2026-10-24}
\bibitem{BrainFramework} Y.~Shmalo, \emph{Brain Framework}, public software repository, 2026.
\url{https://github.com/yspennstate/brain-framework-public/tree/46b935d26d5a2893f7790b60466387dd8da47ef5}
\end{thebibliography}
\end{document}
